\documentclass[10pt,twocolumn]{article}
\usepackage[T1]{fontenc}\usepackage{lmodern,amsmath,amssymb,amsthm,booktabs,array,longtable,geometry,microtype}
\geometry{margin=.77in,columnsep=.26in}
\setlength{\emergencystretch}{2em}
\newtheorem{theorem}{Theorem}
\newtheorem{proposition}[theorem]{Proposition}
\newcommand{\E}{\mathbb E}\newcommand{\dd}{\,d}
\newcommand{\idea}{\par\noindent\emph{Proof idea.} }
\newcommand{\src}[1]{\textnormal{\small [#1]}}
\title{Why a learned identity map can lose accuracy off the data cone}
\author{Theory section draft for the ICML manuscript}
\date{}
\begin{document}\maketitle
A probe can first benefit from learning an identity map and later lose
that benefit while the network continues to learn the training task.
In the small-noise family below, this change of direction occurs after
both clusters meet explicit learning thresholds. The two clusters
already push the probe in opposite directions before the reversal.
What changes is their relative strength. Resolving these forces also
explains why the corresponding linear population model behaves differently.

\section{Model and learned windows}\label{p1thv1:model}
In units of the larger cluster mean \(\mu_1\), draw \(X\) with equal
probability from \(P_1=\mathcal N(e_1,q^2I_2)\) and
\(P_2=\mathcal N(\rho e_2,q^2I_2)\), where
\(q=\sigma/\mu_1,\ \rho=\mu_2/\mu_1,\ \lambda=\rho^2\).
Train the bias-free untied two-layer ReLU population by gradient flow
on \(\tfrac12\E\|f(X)-X\|^2\). Labels \(\beta\) have uniform
probability on the circle; their input and output vectors start equal
to \(\sqrt{S_0}(\cos\beta,\sin\beta)\) and evolve separately. Balance
preserves their common length \(\sqrt m\). Writing
\(a_\theta=(\cos\theta,\sin\theta)\), \(b_\psi=(\cos\psi,\sin\psi)\)
and \(d\mu=m\,d\beta/(2\pi)\), the output is
\[
 f(X)=\int b_\psi(a_\theta^\top X)_+\,d\mu .
\]
For a unit probe \(\xi\), write \(E_\xi=\|\xi-f(\xi)\|^2/2\).
Split its instantaneous derivative into \(D_1^\tau+D_2^\tau\)
by taking the two cluster contributions to the parameter velocity,
each with its mixture weight \(1/2\). A negative rate helps the probe.
At a ReLU gate use the strict-positive selector. Rate signs below
hold for those exact formulas, for classical derivatives where defined,
and for the total time derivative almost everywhere. Define the cluster
response and loss by
\[
 \mathsf m_p=\frac{\E_{P_p}[X_pf_p(X)]}{\E_{P_p}X_p^2},\qquad
 \mathcal L_p=\tfrac12\E_{P_p}\|X-f(X)\|^2 .
\]
\src{RTA sourceconvention, learning; BC channels}

Time on the limiting weak-growth orbit is centered at weak half mass:
\(n=n_\rho(s)=(1+e^{-\lambda s})^{-1}\). For the initialized flow
\[
 \tau=\mu_1^2 t_{\rm phys},\quad s=\tau-\tau_h,\quad
 \tau_h=t_\eta+\lambda^{-1}\log[(1-\eta)/\eta].
\]
Here \(\eta\le1/8\) is the fixed positive entry level and \(t_\eta\)
is its analytically matched auxiliary logistic hit in normalized time.
Finite-noise physical weak mass need not equal \(n\) or hit half at
\(s=0\). Put \(\xi_q=(\sin(q\zeta),-\cos(q\zeta))\). WIT constructs
a fixed positive-width rectangle
\[
 \mathcal P_*=[2/3\pm\delta_\rho]\times[5\pi/9\pm\delta_\zeta]
\]
(where \([a\pm h]=[a-h,a+h]\)), inside
\([13/20,7/10]\times[87/50,7/4]\). On this rectangle all parameters
share \(J=[s_*-r,s_*+r]\), with early and late subintervals
\(J_-=[s_*-r,s_*-r/2]\), \(J_+=[s_*+r/2,s_*+r]\).
The widths, \(r,p_*,\gamma_0>0\) and center \(s_*\) are analytic
orbit constants. \src{WIT orbit, commonwindows, uniformtheorem;
RTA time, rectangle}

\begin{theorem}[Learned reversal and persistence]\label{p1thv1:reversal}
For the initialized population above, \(S_0=q^{10}\), there is
\(q_0>0\), unevaluated, such that for \(0<q\le q_0\) and all
\((\rho,\zeta)\in\mathcal P_*\) the following hold.
On \(J\), \(\mathsf m_1\ge1/2,\ \mathsf m_2\ge1/4\), and
\(\mathcal L_p\le(169/196)\mathcal L_p(0)\) for both clusters, with \(\mathcal L_p(0)\) evaluated at initialization.
The source margins are
\[
 D_1^\tau\le-\tfrac12p_*q^2,\qquad
 D_2^\tau\ge \tfrac12p_*q^2 \quad\text{on }J .
\]
The total rate is at most \(-5\gamma_0q^2/8\) on \(J_-\) and at
least \(5\gamma_0q^2/8\) on \(J_+\). There is an interior minimizer
on \(J\). The errors at \emph{both} endpoints exceed its error by
at least \(\gamma_0rq^2/4\). At each included ratio the probes form an
off-cone angular sector of width \(2\delta_\zeta q\).

There is a common finite \(S>s_*+r\) and, as \(q\downarrow0\),
\(T_q\to\infty,\ \varepsilon_q\to0\) such that uniformly on the
same rectangle, for \(S\le s\le S+T_q\),
\[
 |E_{\xi_q}(\tau_h+s)-1/2|\le\varepsilon_q q^2 .
\]
\src{RTA theoremA, persistence}
\end{theorem}
\idea One fixed set of weak-ancestry labels provides a clock accurate
enough to select the limiting orbit. A separate estimate controls
probe-active background at the \(q^2\) scale. This allows the initialized
source rates to inherit WIT's strict margins. Comparing continuous
errors at the endpoints of \(J_-\) and \(J_+\) gives the rebound.
After the limiting strong input clears the probe, its scaled error is
zero. Uniform convergence on each fixed later interval yields the
unspecified growing horizon by a diagonal argument.
\src{RTA inventory, ratetransfer, theoremA, persistence}

\section{Which motion changes the probe error?}\label{p1thv1:mechanism}
The training cluster specifies the \emph{source} of a gradient.
A moving neuron's cohort is its \emph{carrier}. A cohort whose training
output enters the residual is a \emph{producer}, even when its own
probe gate is closed. These distinctions matter for cluster 2.

Let \(\mathcal K\) be the fixed strong core transported from entry and
\(S_w\) the fixed selected weak cohort (the bulk weak clock).
Their current masses are \(M=\mu(\mathcal K), N=\mu(S_w)\).
On the learned window \(M=1+O(q^{11/10})\),
\(N=n+O(q^{1/8})\), and \(\mathsf m_2=N+O(q^{1/4})\).
On \(\mathcal K\) write \(\theta=qx,\psi=qy\);
on the selected weak chart the input/output directions are
\((\sin(qu),\cos(qu))\) and \((-\sin(qv),\cos(qv))\).
Let \((X_*,Y_*,U_*,V_*)\) be the strong and weak coordinates on WIT's
selected coherent branch at \((\rho,s)\). The branch has strong mass
one, weak mass \(n\), and the resident limit specified in the appendix.
Define \(K(t)=\phi(t)+t\Phi(t)\) for standard normal density and
distribution \(\phi,\Phi\). Set \(A_*=\zeta-X_*>0,\ Z=X_*-Y_*\),
\(P_x=\Phi(\rho X_*)\). The exact exposure and cross-output lag are
\begin{align*}
 \mathcal A_q&=\int_{\mathcal K}(\zeta-x)d\mu,\\
 \mathcal L_q&=\frac{\E_{P_1}[X_1(f_2(X)-X_2)]}{q(1+q^2)} .
\end{align*}
Here \(\mathcal L_q\) denotes a lag, distinguished by its subscript
from cluster loss \(\mathcal L_p\). The core is strictly probe-active.
\src{CPD scope, help; RTA coregap, learningresponse; WIT orbit}

\begin{theorem}[Help and harm on the learned window]\label{p1thv1:channels}
Under Theorem~\ref{p1thv1:reversal}'s hypotheses, uniformly on \(J\),
\begin{align*}
 \mathcal L_q&=Y_*+nK(U_*)+O(q^{1/20}),\\
 q^{-2}D_1^\tau&=-\mathcal A_q\mathcal L_q/2+O(q^{1/20}),\\
 q^{-2}D_2^\tau&=T_{\rm in}+W_{\rm in}+T_{\rm out}+W_{\rm out}
                         +O(q^{1/20}).
\end{align*}
The limiting lag is bounded below by a positive constant, and the
help error is below its signed margin for small \(q\). The four
leading terms are
\begin{align*}
 T_{\rm in}&=-\lambda Z^2P_x/2,&
 W_{\rm in}&=\lambda nZ(V_*-Y_*)P_x/2,\\
 T_{\rm out}&=\rho A_*K(\rho X_*)/2,&
 W_{\rm out}&=-nT_{\rm out}.
\end{align*}
Here \(T_{\rm in}\) combines target and core production.
The target and selected weak contributions to output rotation sum
to \((1-n)T_{\rm out}>0\). On all of \(J\),
\(T_{\rm in}\le0,\ T_{\rm out}>0,\ W_{\rm out}<0\).
The sign \(W_{\rm in}>0\) is proved on \(J\cap(-\infty,1/1000]\);
on any remaining \(J\cap(1/1000,\infty)\) it is OP-W.
Strict limiting signs transfer on compact sets with a positive
margin. \src{BC twoway; CPD table, help; PHS theorem, transfer}
\end{theorem}
\idea In each exact output derivative, retain radial change, input
rotation and output rotation. Split the training residual into target
and fixed-label output producers \emph{before} integrating against
the actual probe weights. For source 1 the target and core contributions
to input rotation and radial change cancel to leading order.
The surviving output motion measures the full signed Gaussian
regression lag. For source 2 both rotations survive. In particular
the leading output motion is positive before weak saturation, even
if the selected weak carriers never fire on the probe.
\src{BC twoway; CPD table, help}

\begin{theorem}[Linear control]\label{p1thv1:linear}
Replace the ReLU by the identity activation with identity targets.
Let the data second-moment matrix be \(A=\operatorname{diag}(a_r)\),
\(a_r>0\), in an orthonormal basis. Suppose input and output weights
are aligned and equal initially almost everywhere and the initial
input Gram matrix in that basis is \(\operatorname{diag}(c_r)\)
with \(0<c_r<1\). Then the global unique linear map stays diagonal
with entries \(\ell_r=[1+(c_r^{-1}-1)e^{-2a_r\tau}]^{-1}\).
For every nonzero \(\xi\) at every finite time,
\[
 E_\xi'=-2\sum_r a_r\ell_r(1-\ell_r)^2\xi_r^2<0,
 \qquad E_{\xi}=E_{-\xi}.
\]
If each mixture-weighted cluster moment \(A_p\) is also diagonal
in this basis, its source derivative has the same sum with
\(a_r\) replaced by \((A_p)_{rr}\), hence is nonpositive; it is
negative if \(\xi\) has support on a positive cluster entry.
For the two Gaussian clusters, aligned isotropic continuum
initialization gives \(G_0=S_0I/2\), \(0<S_0<2\).
The finite-width SIM conclusion requires an exactly positive diagonal
Gram with entries below one on the stated population data.
\src{LP linear-theorem, linear-SIM}
\end{theorem}
\idea Alignment and common diagonal axes preserve independent scalar
logistic modes. Each reduces every squared error component on which it
acts. By contrast, uniform angles in the positive quadrant give
\(G_{11}=G_{22}=S_0/2,\ G_{12}=S_0/\pi\). This Gram leaves the
commuting class when the data eigenvalues differ. The theorem then
gives no sign prediction. Isotropic finite draws can also leave the
class. \src{LP linear-SIM and following positive-quadrant remark}

\section{Where the small-noise theory applies}\label{p1thv1:regime}
\begin{proposition}[Regime map and critical entry]\label{p1thv1:map}
Let \(S_0=q^\kappa\), \(q\downarrow0\), \(13/20\le\rho\le7/10\).
Set \(\omega=(1-\lambda)/2,\ a=\rho K(\rho a),\
P=\Phi(\rho a),\ d=(1-\lambda P)/2,\ s_{\rm sh}=d/\omega\).
At the nominal strong clock \(T_A=\log(A/S_0)\), \(A>0\) fixed,
the six \textbf{C} (comparison-only) thresholds are
\[
 \underbrace{1/\omega}_{\kappa_c},\
 2/\lambda,\ 2/\omega,\ 2/[\lambda(1-\lambda)],\
 4/\lambda,\ 2/\lambda+2/\omega .
\]
They mark, respectively, interior Q1 reaching the strong scale, bulk Q2
crossing the weak axis, last Q1 chart arrival, bulk Q2 reaching an
interior Q1 angle, last Q2 weak-axis crossing and bulk Q2 strong-chart
arrival. The two \textbf{H} (candidate-conditional) thresholds are
\(2s_{\rm sh}/(1-P)\) and \(8s_{\rm sh}/[3(1-P)]\). They test
small shape and a complementary-mass charge assuming transferred
entry, local gain and controlled radial amplification. The entire
band \(2<\kappa<3\) lies below all eight. Ten lies above all except
the last comparison threshold.

At \(S_0=cq^{1/\omega}\) with \(c>0\) fixed,
\(b=(c/A)^\omega\). In the positive-noise target-only comparison the
right-half limiting mass is \(A/4\), half an aligned atom at \(a\)
from Q4 and half a distributed Q1 component above \(a\).
For offsets on a compact subset of \((a,\infty)\), the
right-half inactive fraction equals
\begin{equation}
 I_{\rm ent}(b,\zeta)=\frac12\left[1-\frac4\pi
 \int_0^{H(\zeta)/b}\frac{du}{(1+u^2)^2}\right].
 \label{p1thv1:fraction}
\end{equation}
Here \(H(w)=w\exp[-\omega\int_w^\infty
 (B(r)^{-1}-(\omega r)^{-1})dr]\) on \(w>a\), with
\(B(r)=[r-\rho K(\rho r)]/2\).
\src{RMS map, join}
\end{proposition}
\idea Compare the angular arrival clocks on their specified sectors
with \(T_A\). The ratio \((S_0/A)^\omega/q\) remains finite and
nonzero precisely on the critical scale up to a fixed prefactor.
At positive noise the cross-cluster gate changes the entry coordinate
from the zero-noise value \(b\tan\beta\) to \(H^{-1}(b\tan\beta)\).
Weighting Q1 by \(A\cos^2\beta\,d\beta/(2\pi)\) and adding Q4 yields
\eqref{p1thv1:fraction}. Prescribing this shape as leading data is
a separate step from producing it by training. \src{RMS join, horizons}

\paragraph{Limitations.}
The deep theorem has an unevaluated \(q_0\), an \(O(q)\)-wide sector
and an unspecified persistence horizon. It does not cover the reported
\(176^\circ\) reversal or the other persistent directions outside its
sector. OP-W is the remaining weak-produced input sign. For the
critical family, OP-C1 asks whether a fixed positive fraction stays
inactive indefinitely; the proved survival guarantee is
\(C_*/(1+T)\) on each finite zero-weak-mass resident interval.
OP-C2 asks for the coupled weak passage, seed, clock, exposure and
exact lag. OP-C3 asks for critical trained entry, including
noise-scale angular or weighted-profile precision and tail matching.
At \(q=1/20,\ S_0=1/5000\), \(20^2<5000<20^3\) gives exponent
placement only. No canonical reversal time, magnitude, attribution
or persistent direction is certified here. Any conditional
computer-assisted local evidence has separate hypotheses and is
not a premise. \src{RTA persistence; PHS transfer; RMS open, paper}

\clearpage\onecolumn\appendix
\section{How the initialized theorem is assembled}\label{p1thv1:architecture}
The proof has two sensitive scales. Long before the learned windows,
the weak seed is a power of \(q\), so an order-one clock error changes
its phase. On the learned windows the probe signal is of order \(q^2\),
so a small unscaled population error alone cannot give its sign.
The following dependency list follows RTA's \texttt{inventory, handoff}
and roadmap v1.36, Section 9. S2 and S3 close jointly.
\begin{enumerate}
\item \textbf{S1. Entry at strong half mass.} The early comparison
separates positive strong ancestry, interior Q2 and endpoint layers.
Same-field relative transport transfers the strong profile to a moving
overlap; a mass clock then carries it to \(M=1/2\) at \(10\ell+O(1)\).
Fixed-constant Q2 entry and the recruited tail supply the seed
\(H_0\asymp A_q^{-\lambda}q^{10(1-\lambda)}\), transverse energy
\(O(q^3)\), and the secondary profile. Here \(A_q\) is the bounded
positive strong-entry weight factor. Sources: F/P, ARR, TAIL, Q2,
PROFILE, Q2E, OWPT, MOMC, Q2EP, Q2EC, SEP. The fixed-offset SMCP
audit motivates the moving overlap.
\item \textbf{S2. Tail and distance control, stopped form.} A clock
cohort with small transport defect sets the time; a larger weak
cohort is allowed fixed distortion. A distance envelope and capped
radial gain bound the secondary energy and the complement's
integrated field. The signed-gate budget is applied at each stopped
endpoint. Sources: TAILREP, BULK, RCET, LCRC, DECG, GESG, S2CP.
\item \textbf{S3. Strong field and joint passage.} Keep the strong
donors explicitly and regress the selected weak response in the
partial gate. The response ledger caps weak energy; Q3 strip
budgets pay outside contributions. A joint first-exit argument
and a positive response barrier close the passage to the first
small fixed response \(c_b\). Sources: LOC, CLOCK, CFR, CFA, Q3SB,
PISB, LSB, WPT, RB, with S2CP. This supplies a retained core of
radius \(O(q^{1/20})\) and clock defect \(O(q^{1/8}+q^2\ell)\).
\item \textbf{S4. Match and select the weak orbit.} Return of the
actual weak input to the positive side is needed to use the
contraction estimate. Strong-reference reduction, weighted weak
concentration, and the actual clock select the orbit at a fixed small
weak mass with \(O_\eta(q^{1/20})\) error. Sources: WEAK, TAILREP,
SI (interface), WBR, MP, SRR, WCS.
\item \textbf{S5. Continue to the windows and transfer rates.}
Transverse probe-sector energy pays the background on the \(q^2\)
scale. Continue physical masses and first moments on each fixed time
window, separate every core label from the gate, then compare the two
source derivatives separately. Cluster losses use a full-output
\(L^2\) estimate. WIT supplies limiting margins; LP/LPR supply
common late clearing. Sources: CWI, PSE, FWC, RTA.
\item \textbf{S6. Resolve motion and producers.} Exact source-by-motion
bookkeeping retains the second source's output channel. Fixed-label
producer integration and a signed Gaussian regression give
Theorem~\ref{p1thv1:channels}; PHS extends one sign slightly beyond
half mass. Sources: P, BC, CPD, PHS. This interpretation uses S5;
S6 is not a prerequisite for Theorem~\ref{p1thv1:reversal}.
\end{enumerate}
\paragraph{Constant order.} Keep the ratio box, model constants,
duration cap and the outside-budget scale first. In the joint passage
fix the small response level and the residual/energy bounds before
choosing the small-ratio cutoff \(C_0\); its partition thresholds
are chosen afterwards (the energy guard is \(2c_b\)). Keep the
WPT/RB--WCS order and fixed selection level \(\eta\). Freeze WIT's
original rectangle, windows, gaps and strict margins. Then choose the
finite continuation window, enlarged chart radius and field, distance,
background and learning constants. Choose \(q\) last, including the
conditions that normalized source error is below
\(\min(p_*/8,\gamma_0/16)\) and that the two endpoint errors fit
within the rebound margin. For persistence, common clearing is fixed
before the separate diagonal sequence of noise thresholds.
\src{RTA q0, handoff, persistence; PHS transfer}

\paragraph{Comparator for reconstructing the coefficients.}
Writing \((x,y,u,v)=(X_*,Y_*,U_*,V_*)\), the branch used in the four
terms is the solution, on \(n=(1+e^{-\lambda s})^{-1}\), of
\[
\begin{aligned}
 2x'&=\lambda[nv-x+(1-n)y]\Phi(\rho x),&
 2y'&=-y-nK(u)+\rho(1-n)K(\rho x),\\
 2u'&=(1-n)(\phi(u)-\lambda u)-(nu+y)\Phi(u),&
 2v'&=\rho K(\rho x)-\lambda v .
\end{aligned}
\]
Its incoming state is \((a,a,u_*,a/\lambda)\), where
\(a=\rho K(\rho a)\) and \(\phi(u_*)-a\Phi(u_*)-\lambda u_*=0\).
This identifies the comparator and its phase without using a measured
crossover time. \src{WIT orbit; RTA time}

\section{Reading the experimental ledgers}\label{p1thv1:translation}
The translations use the setup reported for \emph{Swing-by-TwoLayerReLU}
in roadmap v1.36. An exact link identifies a conversion or identity.
A qualitative link compares observations with a theorem on a different
family. A prediction proposes a check. Unreconciled means the two sign
assignments have yet to be matched.
\small\begin{longtable}{@{}>{\raggedright\arraybackslash}p{.18\textwidth}
 >{\raggedright\arraybackslash}p{.42\textwidth}
 >{\raggedright\arraybackslash}p{.32\textwidth}@{}}
\toprule Measured quantity & Theory coordinate & Link and qualification\\\midrule\endhead
Training setup & Finite sum \(\sum_{i=1}^{200}w_i(u_i^\top x)_+\),
identity targets and full-batch loss \(\sum_i\|f(x_i)-x_i\|^2/(2N_{\rm data})\).
Theory uses the corresponding population expectation. &
Exact architecture and loss convention. Finite samples and width
are separate from the continuum theorem.\\
Means and noise & \(\mu_1=3,\ \mu_2=2,\ \sigma=3/20\), giving
\(\rho=2/3,\ q=1/20,\lambda=4/9\). &
Exact change of data units. The cluster-loss ratios are invariant
under the common squared data-unit conversion.\\
Initialization & For aligned units of norm \(\epsilon\),
\(S_0=h\epsilon^2=1/5000\). &
Exact mass definition. A finite isotropic angular draw need not have
the population Gram.\\
Probe and clock & \(\arg\xi=-\pi/2+q\zeta\).
At \(-85^\circ\), \(q\zeta=\pi/36\), hence \(\zeta=5\pi/9\).
\(\tau=9t_{\rm phys}\), \(D_p^{\rm phys}=9D_p^\tau\);
the physical window is \((\tau_h+J)/9\). &
Exact unit conversions. An empirical response-half hit is a different
event from the analytic phase match. \src{RTA time}\\
Learning traces & \(\mathsf m_p\) is a regression coefficient,
\(M,N\) are fixed-cohort radial masses, \(n\) is a logistic comparison.
Match the logged regression numerator and denominator. &
Exact definitions; qualitative proxies outside the deep window.
On the deep learned window \(\mathsf m_2=N+O(q^{1/4})\) and \(N=n+O(q^{1/8})\) hold. \src{RTA learning}\\
Cluster-sourced rate & Contract the full probe residual with the
output derivative generated by one cluster, including its mixture
weight. The source derivatives add to the total.
& Exact for matched definitions. Qualitative for the reported
negative-to-positive switch at canonical initialization.\\
Help and exposure & \(Q_1=1+q^2,\ z=a_\theta^\top\xi,\ 
d_\xi=-(\xi-f(\xi))^\top b_\psi^\perp\).
With current radial averages on the active core,
\(B_0=M Q_1\bar z\,\bar d_\xi\,r_{21}/2\), \(r_{21}=q\mathcal L_q\).
The exact source is \(-B_0-\epsilon_B+r_1\), with geometry defect \(\epsilon_B\) and remaining motions \(r_1\).
Multiply by \(\mu_1^2\) for physical rates. &
Exact product term plus signed remainders. The pilot's near-one
ratio is qualitative evidence about remainder size. Absolute and
signed projections agree only with the matching sign. \src{CPD help}\\
Motion ledger & The exact remainder \(r_2\) after subtracting core input
rotation contains core output rotation, radial motion and noncore
source terms. &
Exact bookkeeping for matched cohorts. Equating the entire remainder
or a reported percentage to one leading term is a prediction.
\src{BC twoway; CPD table}\\
Producer ledger & Split \(X-f(X)\) by the producer of each training
output while retaining the same carrier gate and probe weight.
An inactive weak carrier can shape the strong carrier's velocity. &
Exact distinction. Moving masks, counts and mean-angle evaluation
can change the attribution. Early weak input sign is unreconciled.\\
Architecture and duration & E1/E2 compare architectures; E3 also changes
width or angular initialization. E6 reports persistence at the
canonical probe through physical time 100. &
Exact linear contrast only with LP's hypotheses. Qualitative for finite
canonical observations and duration; no link for directions outside
the proved sector.\\
\bottomrule\end{longtable}\normalsize

\section{Checks for existing logs}\label{p1thv1:checks}
Match source units, carrier cohorts, gates and signed residual producers
before comparing traces. Use the logs' derivative/discretization precision
for exact identities. For extrapolations, fix the tolerance from the
logging protocol before comparing results. A failure at the canonical
pair tests the extrapolation, not a continuum theorem with unevaluated
\(q_0\).

\begin{enumerate}
\item \textbf{Output rotation at crossover.}
Theorem-regime: with \(O_{2,\mathcal K}\) the source-2 strong-core
output-rotation contribution and \(d_{2,*}\) the four-term sum,
\[
 \frac{O_{2,\mathcal K}^{\rm phys}}{\mu_1^2q^2T_{\rm out}}
       =1-n+O(q^{1/20}),\qquad
 \frac{O_{2,\mathcal K}}{D_2}
       =\frac{(1-n)T_{\rm out}}{d_{2,*}}+O(q^{1/20}).
\]
On \(J\), \(T_{\rm out}\) has the positive core gate margin, and WIT
bounds \(d_{2,*}\) below. Prediction outside the theorem regime (E2):
at the measured crossover, separate core output motion from radial
and noncore motion in \(r_2\), and compare the coefficient-normalized
output term with the remaining weak mass deficit. For a share divide
by cluster-2 harm, not by total rate at its zero. A resolved wrong
sign, or systematic failure to track the deficit after normalization,
rejects this extrapolation. The reported post-saturation decay of
\(r_2\) is qualitative evidence. \src{BC twoway; CPD table; WIT uniformtheorem}

\item \textbf{Help product.}
Theorem-regime: \(-D_1^\tau/(q^2\mathcal A_q\mathcal L_q/2)
=1+O(q^{1/20})\) since the scaled denominator has a positive margin.
Prediction outside the theorem regime: compute \(B_0\) with the
same radial weights and full \(X_1\)-weighted lag; distinguish
core output help from the full cluster source. A discrepancy beyond
recording precision after restoring the signed remainder terms
falsifies the bookkeeping match. A wrong-sign lag or a product
residual comparable with help rejects product dominance there.
\src{CPD help, table}

\item \textbf{Linear monotonicity.}
Theorem-regime (exact): verify alignment and the common eigenbasis
of initial Gram and data moments, with \(0<c_r<1\). Every nonzero
probe strictly descends and antipodes agree. A resolved increase or
antipodal difference in that exact flow contradicts the prediction;
a positive cluster source contradicts it when that cluster's moment
is also diagonal. For full-batch sample runs check the actual data
moment. Agreement of E1/E2 finite runs is qualitative without
exact commutation. For narrow or quadrant-initialized E3 runs,
inspect the off-diagonal Gram; leaving the commuting class by itself
predicts no transient. \src{LP linear-theorem, linear-SIM}

\item \textbf{Inactive entry fraction.}
Comparison theorem, prediction outside the trained theorem regime:
choose \(A\) in advance, read the normalized-time slice
\(T_A=\log(A/S_0)\), and substitute \(b=(S_0/A)^\omega/q\) in
\eqref{p1thv1:fraction}. Increasing \(S_0\) increases the predicted
inactive fraction. Use radial mass, initial right-half ancestry and
current activity at that time. A reproducible opposite ordering, or
failure to approach the predicted curve in matched comparisons,
rejects its use as a trained-entry proxy. Changing \(A\) changes
both the fraction and its time slice. A bridge share at crossover
measures a different state. \src{RMS join}

\item \textbf{Position of the sweep.}
At \(\rho=2/3\), \(\kappa_c=18/5\); the reported range of about
\(2.1\) to \(5.9\) [eval] straddles it. The exact horizontal
coordinate is \(\log b=(1-\kappa\omega)\log(1/q)-\omega\log A\).
Recomputing \(\log(1/S_0)/\log(1/q)\) from logged initial norms
and finding all values on the same side would falsify the stated
sweep placement. Prediction outside the theorem regime: larger
powers at fixed \(q,A\) correspond to smaller \(b\), hence smaller
predicted inactive entry mass. An opposite matched entry trend
rejects this extension. The placement gives no finite-\(q\) cutoff
for reversal incidence. \src{RMS map, join}

\item \textbf{Early weak input attribution.}
Unreconciled; prediction outside the theorem regime. The reported
early weak input attribution is negative, whereas the deep
\(W_{\rm in}\) is positive on \(J\cap(-\infty,1/1000]\).
Check whether the observation is even on the corresponding learned
window. Match the weak mass phase and split the weak training output
inside the residual. Keep the strong carrier's input-rotation
channel and its actual rate weights. A negative like-for-like
canonical term rejects the sign extrapolation; an unmatched time
or producer mask leaves the check open. This comparison supplies
no sign theorem for OP-W. \src{CPD table; PHS transfer}
\end{enumerate}
\paragraph{Source keys.}
The suffixes cited above have prefixes RTA \texttt{an07rtav1:},
BC \texttt{an01bcv1:}, CPD \texttt{an01cpdv1:}, PHS \texttt{an01phsv1:},
RMS \texttt{an04rmsv1:}, WIT \texttt{inc:AN06:v2:},
LP \texttt{inc:LP:v1:}. Experimental links follow the roadmap v1.36
experiment interface. Empirical measurements enter no theorem premise.
\end{document}
